\section{正则映射与正则函数}\label{sec:quasiprojective-regular-maps-and-functions}

拟射影簇的概念同时推广了仿射簇和射影簇. 因此, 我们想建立拟射影簇之间的映射概念, 使它在这两种情形下分别专化为已有的多项式映射 (仿射情形) 与正则映射 (射影情形). 事实上, 正确的定义可从射影情形原样沿用, 因此继续使用同一名称.

\begin{Definition}{拟射影簇之间的正则映射}\leavevmode\label{def:regular-map-of-quasiprojective-varieties}
设 $X\subseteq\mathbb{P}^{m}$ 和 $Y\subseteq\mathbb{P}^{n}$ 为拟射影簇. 若对每个 $p\in X$, 都存在同次数的齐次多项式 $f_{0},\ldots,f_{n}\in K[x_{0},\ldots,x_{m}]$, 使得 $p\notin\mathcal{V}(f_{0},\ldots,f_{n})$ 且
\[
  F(a)=[f_{0}(a):\cdots:f_{n}(a)]
  \qquad\text{对所有 }a\in X\setminus\mathcal{V}(f_{0},\ldots,f_{n}),
\]
就称映射 $F:X\to Y$ 为\emph{正则映射}.
\end{Definition}

\BookPairedNote{
特别地, 不仅每个射影簇都可以视为拟射影簇, 而且对一对射影簇, 无论将它们看成射影簇还是拟射影簇, 二者之间的保结构映射都完全相同.
}{
用范畴论的语言说, 这意味着射影簇范畴嵌入拟射影簇范畴.
}
在仿射情形下建立同样的对应要困难得多, 我们将其留到下一节. 也就是说, 要证明一对仿射簇之间的多项式映射与它们典范地对应的拟射影簇之间的正则映射完全相同. 现在只考察这一现象的一个具体例子.

\begin{Example}{仿射开片之间的正则映射}\leavevmode\label{ex:regular-map-of-affine-patches}
考察拟射影簇
\[
  \begin{aligned}
  X&=\{[a_{0}:a_{1}:a_{2}]\in\mathbb{P}^{2}\mid a_{0}\ne0\},\\
  Y&=\{[b_{0}:b_{1}]\in\mathbb{P}^{1}\mid b_{0}\ne0\},
  \end{aligned}
\]
它们分别是仿射开片 $\mathbb{A}^{2}_{0}$ 和 $\mathbb{A}^{1}_{0}$. 此时
\[
  F([a_{0}:a_{1}:a_{2}])=[a_{0}^{2}:a_{0}a_{2}+a_{1}^{2}]
\]
给出正则映射 $F:X\to Y$, 因为其坐标函数是多项式. 将 $Y$ 与 $\mathbb{A}^{1}$ 等同后, 这是映射
\[
  [a_{0}:a_{1}:a_{2}]\mapsto\frac{a_{0}a_{2}+a_{1}^{2}}{a_{0}^{2}},
\]
看似是有理的, 但不是多项式的. 不过, 再将 $X$ 与 $\mathbb{A}^{2}$ 等同, 就变成映射 $(a_{1},a_{2})\mapsto a_{2}+a_{1}^{2}$, 的确是仿射簇 $\mathbb{A}^{2}$ 和 $\mathbb{A}^{1}$ 之间的多项式映射.
\end{Example}

证明仿射簇之间每个正则映射都是多项式映射的困难在于, 当把一对仿射簇视为拟射影簇时, 它们之间的正则映射可能只能局部用多项式描述, 而仿射簇之间的多项式映射必须在整体上由多项式给出. 对仿射簇而言, 局部的多项式表示蕴含整体的多项式表示, 这正是下一节 Theorem~\ref{thm:regular-functions-on-affine-varieties-are-polynomial} 的内容.

拟射影簇之间的正则映射有一个重要特例: 目标为拟射影簇 $\mathbb{A}^{1}=K$. 可以把这样的映射看成拟射影簇上的\emph{正则函数}. 的确, 即使主要关心的是射影簇, 这也给出了引入拟射影簇的一个重要动机: 直到现在, 我们还没有射影情形中正则函数的概念, 但现在可将它定义为正则映射的特例.

更具体地说, 考察 Definition~\ref{def:regular-map-of-quasiprojective-varieties} 在 $Y=j_{0}(\mathbb{A}^{1})$ 时的特化. 展开通过
\[
  [a_{0}:a_{1}]\in j_{0}(\mathbb{A}^{1})\longleftrightarrow\frac{a_{1}}{a_{0}}\in K
\]
将 $j_{0}(\mathbb{A}^{1})$ 与 $\mathbb{A}^{1}=K$ 等同的过程, 就得到以下定义.

\begin{Definition}{拟射影簇上的正则函数}\leavevmode\label{def:regular-function-on-quasiprojective-variety}
设 $X\subseteq\mathbb{P}^{m}$ 为拟射影簇. 若对每个 $p\in X$, 都存在同次数的齐次多项式 $f,g\in K[x_{0},\ldots,x_{m}]$, 使得 $p\notin\mathcal{V}(g)$ 且
\[
  F(a)=\frac{f(a)}{g(a)}\qquad\text{对所有 }a\in X\setminus\mathcal{V}(g),
\]
就称函数 $F:X\to K$ 为\emph{正则函数}.
\end{Definition}

换句话说, 正则函数 $F:X\to K$ 是可以局部用有理函数 (不一定是多项式) 描述的函数. 下一例给出一个描述在整体上有效的正则函数.

\begin{Example}{$\mathbb{P}^{1}\setminus\mathcal{V}(xy)$ 上的正则函数}\leavevmode\label{ex:regular-function-on-twice-punctured-projective-line}
考察拟射影簇 $X=\mathbb{P}^{1}\setminus\mathcal{V}(xy)$. 注意, $X$ 是 $\mathbb{P}^{1}$ 中两点 $[1:0]$ 和 $[0:1]$ 的补集. 那么
\[
  F([a:b])=\frac{a^{2}+b^{2}}{ab}
\]
是 $X$ 上正则函数的一个例子, 原因仅在于它是两个同次数多项式的商, 且分母在 $X$ 上从不消没.
\end{Example}

正则函数的和, 积及标量倍数仍正则 (Exercise~\ref{ex:regular-functions-form-algebra}), 因此拟射影簇上的正则函数的集合自然构成一个 $K$-代数. 我们用仿射情形中表示坐标环的同一记号表示这个 $K$-代数.

\begin{Definition}{正则函数环}\leavevmode\label{def:ring-of-regular-functions}
设 $X$ 为拟射影簇. $X$ 上的\emph{正则函数环}记为 $K[X]$, 是所有正则函数 $F:X\to K$ 构成的 $K$-代数.
\end{Definition}

提醒读者, $K[X]$ 对仿射簇现在具有双重含义: 可能是它的坐标环, 也可能是它的正则函数环. 下一节将通过证明这两个环同构来解决这一问题.

和正则映射的情形一样, 主要的微妙之处在于 Definition~\ref{def:regular-function-on-quasiprojective-variety} 是局部的. 这容许正则函数 $F$ 在不同点 $p\in X$ 附近由不同的多项式商 $f/g$ 表示. 不过, 事实上有一类很好的拟射影簇, 在其中不会出现这一局部与整体的微妙差别, 可以用在整体上定义的商显式地描述正则函数环.

\begin{Lemma}{$\mathbb{P}^{m}$ 的开子集上的正则函数}\leavevmode\label{lem:regular-functions-on-projective-open-subsets-are-global-quotients}
设 $X\subseteq\mathbb{P}^{m}$ 为开集. 给定任意正则函数 $F:X\to K$, 存在同次数的齐次多项式 $f,g\in K[x_{0},\ldots,x_{m}]$, 使 $X\cap\mathcal{V}(g)=\varnothing$, 且
\[
  F(a)=\frac{f(a)}{g(a)}\qquad\text{对所有 }a\in X.
\]
\end{Lemma}
\begin{Proof}
令 $p\in X$. 由正则函数的定义, 存在同次数的齐次多项式 $f,g\in K[x_{0},\ldots,x_{m}]$, 使 $g(p)\ne0$ 且
\[
  F(a)=\frac{f(a)}{g(a)}\qquad\text{对所有 }a\in X\setminus\mathcal{V}(g).
\]
必要时约去因子后 (见 Exercise~\ref{ex:cancel-common-factors-in-regular-function}), 可假设多项式 $f$ 和 $g$ 没有公共因子. 我们要证明 $\mathcal{V}(g)\cap X=\varnothing$, 从而 $F$ 的上述商表示在整个 $X$ 上有效. 为此, 只须证明 $g$ 没有在 $X$ 的一点消没的不可约因子.

为导出矛盾, 假设 $h$ 是 $g$ 的一个在点 $q\in X$ 消没的不可约因子. 由正则函数的定义, 存在同次数的齐次多项式 $f_{q},g_{q}\in K[x_{0},\ldots,x_{m}]$, 使 $g_{q}(q)\ne0$ 且
\[
  F(a)=\frac{f_{q}(a)}{g_{q}(a)}\qquad\text{对所有 }a\in X\setminus\mathcal{V}(g_{q}).
\]
注意, $X\setminus\mathcal{V}(g)$ 和 $X\setminus\mathcal{V}(g_{q})$ 都是 $\mathbb{P}^{m}$ 中的非空开集, 因此由 Exercise~\ref{ex:irreducibility-criterion-by-nonempty-open-sets}, 它们的交 $X\setminus(\mathcal{V}(g)\cup\mathcal{V}(g_{q}))$ 也是 $\mathbb{P}^{m}$ 中的非空开集. 此外, 对任意 $a\in X\setminus(\mathcal{V}(g)\cup\mathcal{V}(g_{q}))$, 有
\[
  \frac{f(a)}{g(a)}=\frac{f_{q}(a)}{g_{q}(a)}.
\]
所以 $fg_{q}-f_{q}g\in\mathcal{I}(X\setminus(\mathcal{V}(g)\cup\mathcal{V}(g_{q})))$. 由于 $X\setminus(\mathcal{V}(g)\cup\mathcal{V}(g_{q}))$ 在 $\mathbb{P}^{m}$ 中非空且开, 由 Proposition~\ref{prop:nonempty-open-in-irreducible-space-is-dense} 知它稠密, 再由 Example~\ref{ex:projective-closure-of-arbitrary-subset} 得
\[
  \mathcal{I}(X\setminus(\mathcal{V}(g)\cup\mathcal{V}(g_{q})))
    =\mathcal{I}(\mathbb{P}^{m})=\{0\}.
\]
因此 $fg_{q}-f_{q}g=0$. 按假设 $g(q)=0$, 而 $g_{q}(q)\ne0$, 故必有 $f(q)=0$.

到这里已证明 $h(q)=0$ 蕴含 $f(q)=0$, 因而 $\mathcal{V}(h)\cap X\subseteq\mathcal{V}(f)$. 由于 $\mathcal{V}(h)\cap X$ 是不可约簇 $\mathcal{V}(h)$ 的非空开集, 它稠密, 从而 $\mathcal{V}(h)\subseteq\mathcal{V}(f)$. 由于假设 $h$ 不可约, 零点定理蕴含 $h$ 整除 $f$, 与 $f$ 和 $g$ 没有公共因子的假设矛盾. 由此矛盾可知, 表达式 $F=f/g$ 就是所需的整体上定义的商.
\end{Proof}

现在考察上述引理的几个具体应用.

\begin{Example}{$\mathbb{P}^{m}$ 的正则函数环}\leavevmode\label{ex:regular-functions-on-projective-space-are-constant}
考察 $X=\mathbb{P}^{m}$. 由 Lemma~\ref{lem:regular-functions-on-projective-open-subsets-are-global-quotients}, $X$ 上的正则函数 $F$ 可表示为 $F=f/g$, 其中 $g$ 在 $\mathbb{P}^{m}$ 上处处不消没. 这只有在 $g$ 是常数时才可能. 又因为 $f$ 和 $g$ 同次数且齐次, 所以 $f$ 也是常数. 换句话说, $\mathbb{P}^{m}$ 上的每个正则函数都是常值函数, 因此
\[
  K[\mathbb{P}^{m}]=K.
\]
事实上, 这是一个普遍现象的特例: 任何射影簇上都没有非常值的正则函数. 不过, 一般情形的证明要困难得多, 要到本章末尾才能完成.
\end{Example}

\begin{Example}{$\mathbb{P}^{1}\setminus\mathcal{V}(xy)$ 上的正则函数环}\leavevmode\label{ex:regular-function-ring-of-twice-punctured-projective-line}
如 Example~\ref{ex:regular-function-on-twice-punctured-projective-line} 一样, 令
\[
  X=\mathbb{P}^{1}\setminus\mathcal{V}(xy).
\]
令 $F$ 为 $X$ 上的正则函数. 由 Lemma~\ref{lem:regular-functions-on-projective-open-subsets-are-global-quotients}, 可表示 $F=f/g$, 其中 $f,g$ 是同次数的齐次多项式, 且 $g$ 只有在 $x=0$ 或 $y=0$ 时才消没. 这意味着 $g$ 形如 $x^{j}y^{k}$, 其中 $j,k\ge0$. 由此, 在 $f/g$ 中取 $y=1$, 就得到关于 $x$ 和 $1/x$ 的多项式. 例如, 对 Example~\ref{ex:regular-function-on-twice-punctured-projective-line} 中的 $F$, 有
\[
  F(x,1)=\frac{x^{2}+1}{x}=x+\frac{1}{x}\in K[x,x^{-1}].
\]
\BookPairedNote{
这一去齐次化过程给出一个函数 $K[X]\to K[x,x^{-1}]$. Exercise~\ref{ex:isomorphism-with-laurent-polynomial-ring} 建议读者证明此函数是双射同态, 从而证明
\[
  K[X]\cong K[x,x^{-1}].
\]
}{
回顾 $K[x,x^{-1}]$ 表示 $K(x)$ 中由 $x$ 和 $x^{-1}$ 生成的子代数.
}
\end{Example}

本节最后简要讨论拟射影情形下的同构. 定义大概正如读者所料.

\begin{Definition}{拟射影簇之间的同构}\leavevmode\label{def:isomorphism-of-quasiprojective-varieties}
拟射影簇 $X\subseteq\mathbb{P}^{m}$ 和 $Y\subseteq\mathbb{P}^{n}$ 之间的\emph{同构}, 是具有正则逆映射 $G:Y\to X$ 的正则映射 $F:X\to Y$. 若 $X$ 与 $Y$ 之间存在一个同构, 就称 $X$ 与 $Y$ \emph{同构}, 并记 $X\cong Y$.
\end{Definition}

\begin{Example}{拟射影簇之间的一个同构}\leavevmode\label{ex:isomorphism-of-punctured-projective-line-and-conic}
如 Examples~\ref{ex:regular-function-on-twice-punctured-projective-line} 和~\ref{ex:regular-function-ring-of-twice-punctured-projective-line} 一样, 考察 $X=\mathbb{P}^{1}\setminus\mathcal{V}(xy)$, 并令 $Y$ 为拟射影簇
\[
  Y=\mathcal{V}(w^{2}-xy)\setminus\mathcal{V}(w)\subseteq\mathbb{P}^{2}.
\]
则由
\[
  F([a:b])=[ab:a^{2}:b^{2}]
\]
给出的正则映射 $F:X\to Y$ 是同构, 因为 $G([a:b:c])=[a:c]$ 是它的正则逆映射. 建议读者画出该同构在 $\mathbb{R}$ 上的图形.
\end{Example}

正如所希望的, 正则函数环在同构下保持. 为了证明这一点, 我们需要仿射簇理论中已经遇到过的概念在拟射影情形下的对应: 拉回同态.

\begin{Definition}{拉回同态}\leavevmode\label{def:pullback-on-regular-function-rings}
令 $X,Y$ 为拟射影簇, $F:X\to Y$ 为正则映射. $F$ 诱导的\emph{拉回同态}为
\[
  \begin{aligned}
    F^{*}:K[Y]&\to K[X]\\
    F^{*}(G)&=G\circ F.
  \end{aligned}
\]
\end{Definition}

读者应该停下来确认, $F^{*}$ 的确将 $Y$ 上的正则函数变为 $X$ 上的正则函数 (Exercise~\ref{ex:pullback-regular-functions-is-algebra-homomorphism}). 此外, 第~\ref{sec:pullback-homomorphisms}~节描述的仿射情形下拉回的基本性质可原样沿用到拟射影情形. 因而, 和仿射情形相同的证明给出如下结论: 正则函数环在同构下不变.

\begin{Proposition}{$X\cong Y$ 蕴含 \ensuremath{K[X]\cong K[Y]}}\leavevmode\label{prop:quasiprojective-isomorphism-induces-regular-function-ring-isomorphism}
若 $F:X\to Y$ 是拟射影簇之间的同构, 则拉回同态 $F^{*}:K[Y]\to K[X]$ 是 $K$-代数的同构.
\end{Proposition}
\begin{Proof}
Exercise~\ref{ex:ring-invariance-under-quasiprojective-isomorphism}.
\end{Proof}

\BookPairedNoteLeft{
不同于仿射情形, Proposition~\ref{prop:quasiprojective-isomorphism-induces-regular-function-ring-isomorphism} 的逆命题不成立. 如 Example~\ref{ex:regular-functions-on-projective-space-are-constant} 所见, 对任意 $m$ 都有 $K[\mathbb{P}^{m}]=K$, 然而当 $m\ne n$ 时 $\mathbb{P}^{m}\not\cong\mathbb{P}^{n}$.
}{
现在已经熟悉拟射影簇及其间保持结构的正则映射, 可以回到仿射情形, 把仿射簇和多项式映射的理论确切地置于拟射影簇和正则映射的框架之中. 这是下一节的主题.
}

\BookExerciseHeading
\begin{BookExercises}
  \item\label{ex:quasiprojective-regular-composition-and-isomorphism-relation}
  \begin{enumerate}
    \item 证明正则映射的复合仍是正则映射.
    \item 证明同构是拟射影簇的集合上的等价关系.
  \end{enumerate}
  \item\label{ex:regular-functions-form-algebra} 设 $X\subseteq\mathbb{P}^{n}$ 为拟射影簇, $F,G:X\to K$ 为正则函数, $\lambda\in K$. 证明和 $F+G:X\to K$, 积 $FG:X\to K$, 以及标量倍数 $\lambda F:X\to K$ 都是正则函数, 并验证这些运算满足 $K$-代数的公理.
  \item\label{ex:regular-function-ring-is-reduced} 证明拟射影簇的正则函数环是既约的.
  \item\label{ex:cancel-common-factors-in-regular-function} 设 $X\subseteq\mathbb{P}^{n}$ 为不可约拟射影簇, 考察齐次多项式 $f,g,h\in K[x_{0},\ldots,x_{n}]$, 其中 $f$ 和 $g$ 同次数. 令 $F:X\to K$ 为正则函数. 若 $\mathcal{V}(gh)\ne X$, 且
  \[
    F(a)=\frac{f(a)h(a)}{g(a)h(a)}
    \qquad\text{对每个 }a\in X\setminus\mathcal{V}(gh),
  \]
  证明
  \[
    F(a)=\frac{f(a)}{g(a)}\qquad\text{对每个 }a\in X\setminus\mathcal{V}(g).
  \]
  \item\label{ex:isomorphism-with-laurent-polynomial-ring} 令 $X=\mathbb{P}^{1}\setminus\mathcal{V}(xy)$. 证明 $K[X]\cong K[x,x^{-1}]$.
  \item\label{ex:regular-function-ring-of-affine-space} 利用 Lemma~\ref{lem:regular-functions-on-projective-open-subsets-are-global-quotients}, 证明将 $\mathbb{A}^{n}$ 视为仿射开片 $\mathbb{A}^{n}_{0}\subseteq\mathbb{P}^{n}$ 时, 它的正则函数环是 $K[x_{1},\ldots,x_{n}]$.
  \item\label{ex:pullback-regular-functions-is-algebra-homomorphism} 设 $X\subseteq\mathbb{P}^{m}$ 和 $Y\subseteq\mathbb{P}^{n}$ 为拟射影簇, $F:X\to Y$ 为正则映射.
  \begin{enumerate}
    \item 对任意 $G\in K[Y]$, 证明 $G\circ F\in K[X]$.
    \item 证明 $F^{*}:K[Y]\to K[X]$ 是 $K$-代数同态.
  \end{enumerate}
  \item\label{ex:ring-invariance-under-quasiprojective-isomorphism} 验证第~\ref{sec:pullback-homomorphisms}~节中拉回的相关性质可推广到正则映射与正则函数, 由此证明 Proposition~\ref{prop:quasiprojective-isomorphism-induces-regular-function-ring-isomorphism}.
  \item\label{ex:regular-map-from-projective-space-is-polynomial} 设 $Y\subseteq\mathbb{P}^{n}$ 为拟射影簇. 证明每个正则映射 $F:\mathbb{P}^{m}\to Y$ 都是多项式映射.
\end{BookExercises}
